\documentclass{ctexart}
\usepackage{amsmath, amssymb}
\usepackage{geometry}
\geometry{a4paper, margin=2cm}

\title{\textbf{第5章 质点的动量矩与动量矩守恒} --- 例题解答}
\author{}
\date{}

\begin{document}
\maketitle

\begin{enumerate}

\item（例 6.10）质量为 $m$ 的小球，悬挂在一根轻杆下端，杆的上端固定在一个可绕竖直轴旋转的铰链上。设小球在水平面内作半径为 $r$ 的匀速率圆周运动，角速度为 $\omega$，轻杆与竖直方向之间的夹角为 $\alpha$。求小球对 $A$ 点（铰链）和 $B$ 点（杆与小球连接点）的动量矩。

\textbf{解：}小球作匀速圆周运动，速度 $v = r\omega$，动量 $p = mv$。

对 $A$ 点（铰链，位于转轴上）：
动量矩方向沿 $z$ 轴（竖直向上），大小为
\[
|L_A| = mrv = mr^{2}\omega
\]
$L_A$ 的大小和方向都固定不变。

对 $B$ 点（杆与小球连接点）：
$B$ 点与小球在同一水平面，$r_1 = \dfrac{r}{\sin\alpha}$，
\[
|L_B| = r_1 m v = m r_1 r\omega = \frac{mr^{2}\omega}{\sin\alpha}
\]
$L_B$ 的大小不变，但方向不断变化。
$L_B$ 沿 $z$ 轴的投影为
\[
L_z = L_B\sin\alpha = mr^{2}\omega
\]
即 $L_B$ 沿 $z$ 轴的投影不变，且等于小球对 $A$ 点（即对 $z$ 轴）的动量矩。

\item 一质点 $m$，速度为 $v$，$A$、$B$、$C$ 分别为三个参考点，此时 $m$ 相对三个点的距离分别为 $d_1$、$d_2$、$d_3$。求此时刻质点对三个参考点的动量矩。

\textbf{解：}动量矩 $L = r_\perp mv$，其中 $r_\perp$ 是参考点到速度方向直线的垂直距离（垂距）。
\[
L_A = d_1 mv,\quad L_B = d_1 mv,\quad L_C = 0
\]
动量矩只与垂距有关，不是到参考点的直接距离。

\item（例题 1）系在绳子的一端，绳的另一端穿过水平光滑桌面中央的小洞 $O$，起初下面用手拉着不动，质点在桌面上绕 $O$ 作匀速圆周运动。然后慢慢下拉绳子，使它在桌面上那一段缩短。质点绕 $O$ 的角速度 $\omega$ 如何随半径 $r$ 变化？

\textbf{解：}质点受到的是有心力（绳子拉力始终指向 $O$），对 $O$ 点的力矩为零，角动量守恒。
角动量 $L = mrv = mr^{2}\omega$ 为常量，故
\[
\omega \propto \frac{1}{r^{2}}
\]
即角速度反比于半径的平方，半径减小则角速度增大。

\item（例题 2）两个同样重的小孩，各抓着跨过滑轮绳子的两端。一个孩子用力向上爬，另一个则抓住绳子不动。若滑轮的质量和轴上的摩擦都可忽略，哪一个小孩先到达滑轮？又：两个小孩重量不等时情况如何？

\textbf{解：}以滑轮轴为参考点，小孩和滑轮组成系统。
(1) 两小孩等重：系统所受外力矩（两小孩重力的力矩）大小相等、方向相反，合外力矩为零，系统角动量守恒。
起初 $v_1 = v_2 = 0$，$J = 0$；运动后仍 $J = 0$，即 $v_1 = v_2$，故同时到达滑轮。

(2) 两小孩不等重（$m_1 > m_2$）：合外力矩 $M = (m_2 - m_1)gR \neq 0$。
由角动量定理 $\dfrac{\mathrm{d}J}{\mathrm{d}t} = M$，体轻的小孩上升得快，先到达滑轮。

\item（书中例题 6.11，p.214）人造卫星在椭圆轨道上运行，地球中心可看作固定点，近地点离地面的距离为 $439\,\text{km}$，远地点离地面的距离为 $2384\,\text{km}$，近地点速度为 $8.12\,\text{km/s}$，地球半径为 $6370\,\text{km}$。求卫星在远地点的速度。

\textbf{解：}卫星受地球引力为有心力，对地球中心的力矩为零，动量矩守恒。
\[
mv_A R_A = mv_B R_B
\]
\[
v_B = v_A \frac{R_A}{R_B}
\]
其中
\[
R_A = 439 + 6370 = 6809\,\text{km},\quad
R_B = 2384 + 6370 = 8754\,\text{km}
\]
代入得
\[
v_B = 8.12 \times \frac{6809}{8754} \approx 6.32\,\text{km/s}
\]
近地点速度快，远地点速度慢，两点动量不同但动量矩相同。

\item（书中例题 6.12，p.215，重点）质量为 $m$ 的小球系在绳子的一端，绳穿过一铅直套管，使小球限制在一光滑水平面上运动。先使小球以速度 $v_0$ 绕管心作半径为 $r_0$ 的圆周运动，然后向下拉绳，使小球轨迹最后成为半径为 $r$ 的圆。试求：
\begin{enumerate}
  \item 小球距管心 $r$ 时速度 $v$ 的大小；
  \item 绳从 $r_0$ 缩短到 $r$ 过程中，力 $F$ 所作的功。
\end{enumerate}

\textbf{解：}(1) 绳子拉力为有心力，对 $O$ 点力矩为零，动量矩守恒：
\[
mv_0 r_0 = mvr \;\Longrightarrow\; v = v_0 \frac{r_0}{r}
\]

(2) 由动能定理，力 $F$ 作的功等于小球动能的增量：
\[
A = \frac12 mv^{2} - \frac12 mv_0^{2}
= \frac12 mv_0^{2}\left(\frac{r_0^{2}}{r^{2}} - 1\right)
\]

\item 发射一宇宙飞船去考察一质量为 $M$、半径为 $R$ 的行星，当飞船静止于空间距行星中心 $4R$ 时，以速度 $v_0$ 发射一质量为 $m$ 的仪器。要使该仪器恰好掠过行星表面，求 $\theta$ 角及着陆滑行的初速度。

\textbf{解：}行星引力为有心力，系统机械能守恒且对行星中心的动量矩守恒。

机械能守恒：
\[
\frac12 mv_0^{2} - G\frac{Mm}{4R} = \frac12 mv^{2} - G\frac{Mm}{R}
\]

动量矩守恒：
\[
mv_0(4R)\sin\theta = mvR
\]

联立解得：
\[
\sin\theta = \frac14\sqrt{1 + \frac{3GM}{2Rv_0^{2}}},\qquad
v = 4v_0\sin\theta = v_0\sqrt{1 + \frac{3GM}{2Rv_0^{2}}}
\]

\end{enumerate}

\end{document}
